% Options for packages loaded elsewhere
\PassOptionsToPackage{unicode}{hyperref}
\PassOptionsToPackage{hyphens}{url}
%
\documentclass[
]{article}
\usepackage[left=0.7in,right=0.7in]{geometry}
\usepackage{amsmath,amssymb}
\usepackage{iftex}
\ifPDFTeX
  \usepackage[T1]{fontenc}
  \usepackage[utf8]{inputenc}
  \usepackage{textcomp} % provide euro and other symbols
\else % if luatex or xetex
  \usepackage{unicode-math} % this also loads fontspec
  \defaultfontfeatures{Scale=MatchLowercase}
  \defaultfontfeatures[\rmfamily]{Ligatures=TeX,Scale=1}
\fi
\usepackage{lmodern}
\ifPDFTeX\else
  % xetex/luatex font selection
\fi
% Use upquote if available, for straight quotes in verbatim environments
\IfFileExists{upquote.sty}{\usepackage{upquote}}{}
\IfFileExists{microtype.sty}{% use microtype if available
  \usepackage[]{microtype}
  \UseMicrotypeSet[protrusion]{basicmath} % disable protrusion for tt fonts
}{}
\makeatletter
\@ifundefined{KOMAClassName}{% if non-KOMA class
  \IfFileExists{parskip.sty}{%
    \usepackage{parskip}
  }{% else
    \setlength{\parindent}{0pt}
    \setlength{\parskip}{6pt plus 2pt minus 1pt}}
}{% if KOMA class
  \KOMAoptions{parskip=half}}
\makeatother
\usepackage{xcolor}
\usepackage{longtable,booktabs,array}
\usepackage{calc} % for calculating minipage widths
% Correct order of tables after \paragraph or \subparagraph
\usepackage{etoolbox}
\makeatletter
\patchcmd\longtable{\par}{\if@noskipsec\mbox{}\fi\par}{}{}
\makeatother
% Allow footnotes in longtable head/foot
\IfFileExists{footnotehyper.sty}{\usepackage{footnotehyper}}{\usepackage{footnote}}
\makesavenoteenv{longtable}
\setlength{\emergencystretch}{3em} % prevent overfull lines
\providecommand{\tightlist}{%
  \setlength{\itemsep}{0pt}\setlength{\parskip}{0pt}}
\setcounter{secnumdepth}{-\maxdimen} % remove section numbering
\ifLuaTeX
  \usepackage{selnolig}  % disable illegal ligatures
\fi
\IfFileExists{bookmark.sty}{\usepackage{bookmark}}{\usepackage{hyperref}}
\IfFileExists{xurl.sty}{\usepackage{xurl}}{} % add URL line breaks if available
\urlstyle{same}
\hypersetup{
  hidelinks,
  pdfcreator={LaTeX via pandoc}}

\author{}
\date{}

\begin{document}

\hypertarget{secure-atm---design-document---group-13}{%
\section{Secure ATM - Final Design Document - Group
13}\label{secure-atm---design-document---group-13}}

\begin{longtable}[]{@{}llll@{}}
\toprule\noalign{}
S. No. & Roll No. & GitHub ID & Name \\
\midrule\noalign{}
\endhead
\bottomrule\noalign{}
\endlastfoot
1 & 22110047 & \texttt{bp0609} & Bhavik Patel \\
2 & 22110089 & \texttt{guntas-13} & Guntas Singh Saran \\
3 & 22110184 & \texttt{Zeenu03} & Jinil Patel \\
4 & 23110276 & \texttt{rishabhh-7} & Rishabh Jogani \\
\end{longtable}

\hypertarget{system-overview}{%
\subsection{1. System Overview}\label{system-overview}}

We build two C++17 programs - \texttt{bank} (server) and \texttt{atm}
(client) - that communicate over TCP through an encrypted, authenticated
channel. Security is bootstrapped from a single \textbf{pre-shared auth
file} containing a 256-bit master secret, distributed out-of-band over a
trusted channel before any communication begins. The bank maintains an
\textbf{in-memory account ledger}; the ATM is stateless and executes one
transaction per invocation.

Because both sides already share a secret, we use \textbf{symmetric
cryptography only} - no TLS, no certificates, no public-key operations -
the is the correct and sufficient choice when a pre-shared key exists.

\begin{center}\rule{0.5\linewidth}{0.5pt}\end{center}

\hypertarget{how-the-protocol-runs}{%
\subsection{2. How the Protocol Runs}\label{how-the-protocol-runs}}

\hypertarget{key-setup}{%
\subsubsection{Key Setup}\label{key-setup}}

On first launch, the bank generates 256 random bits
(\texttt{K\_MASTER}), writes them to the auth file, and prints
\texttt{"created\textbackslash{}n"}. Both sides then derive three
purpose-specific keys using
\href{https://datatracker.ietf.org/doc/html/rfc5869}{HKDF-SHA256}:

\begin{itemize}
\tightlist
\item
  \texttt{K\_ENC} - for
  \href{https://en.wikipedia.org/wiki/Galois/Counter_Mode}{AES-256-GCM}
  encryption
\item
  \texttt{K\_MAC} - for outer HMAC-SHA256 integrity
\item
  \texttt{K\_CARD} - for deterministic card-secret derivation
\end{itemize}

\hypertarget{wire-format}{%
\subsubsection{Wire Format}\label{wire-format}}

Every message is length-prefixed (4 bytes, big-endian) followed by an
encrypted frame:

\begin{verbatim}
Wire Frame:
+----------+-----------+--------------------------------+------------+-------------+
| SeqNo    | Nonce     | AES-256-GCM Ciphertext + Tag   | Timestamp  | HMAC-SHA256 |
| (8 bytes)|(12 bytes) | (variable)                     | (8 bytes)  | (32 bytes)  |
+----------+-----------+--------------------------------+------------+-------------+
                                                                          |
            HMAC covers everything from SeqNo to Timestamp ---------------+
\end{verbatim}

The outer HMAC covers everything from SeqNo through Timestamp, providing
a defence-in-depth authentication layer on top of GCM's built-in tag.

\hypertarget{initial-connection-protocol}{%
\subsubsection{Initial Connection
Protocol}\label{initial-connection-protocol}}

Before any encrypted messages flow, the bank must be running and the ATM
must establish a TCP connection to it.

\begin{verbatim}
BANK (server)                                ATM (client)
    |                                            |
    |  1. Parse CLI flags (-p port, -s authfile) |  1. Parse CLI flags (-i ip, -p port,
    |  2. Generate K_MASTER (256 bits CSPRNG)    |     -s authfile, -c cardfile, ...)
    |  3. Write auth file                        |  2. Validate all inputs (account name,
    |  4. Print "created\n" to stdout            |     amount format, file names, port range)
    |  5. Derive K_ENC, K_MAC, K_CARD via HKDF   |  3. Read auth file -> derive K_ENC,
    |  6. Create TCP socket                      |     K_MAC, K_CARD via HKDF
    |  7. setsockopt(SO_REUSEADDR)               |  4. Read card file or direct card content (if not -n mode)
    |  8. bind(0.0.0.0, port)                    |
    |  9. listen(backlog)                        |
    |  10. Accept loop begins                    |
    |     |                                      |
    |     |<-------- TCP SYN --------------------|  5. connect(bank_ip, bank_port)
    |     |--------- SYN-ACK ------------------> |
    |     |<-------- ACK ----------------------- |
    |     |                                      |
    |     |  accept() returns new socket fd      |  Connection established
    |     |  Set SO_RCVTIMEO = 3 seconds         |  Set SO_RCVTIMEO = 3 seconds
    |     |                                      |
    |     |  ... encrypted protocol begins ...   |
\end{verbatim}

\textbf{Startup ordering matters:} The bank must be fully running
(listening on its port and having printed
\texttt{"created\textbackslash{}n"}) before any ATM instance connects.
If the ATM connects before the bank is ready, \texttt{connect()} fails
and the ATM exits with code 63.

\textbf{One connection per invocation:} Each ATM process opens exactly
one TCP connection, performs one transaction (3 encrypted messages), and
exits. The bank uses a thread-per-client accept loop with a hard
connection cap (\texttt{MAX\_ACTIVE\_CLIENTS = 60}). Every connection is
subject to a 3-second socket timeout.

\newpage
\hypertarget{transaction-flow-3-messages}{%
\subsubsection{Transaction Flow (3
messages)}\label{transaction-flow-3-messages}}

\begin{verbatim}
ATM                                      BANK
 |                                        |
 |---- TCP Connect ---------------------> |
 |                                        |  Generate CHALLENGE R_B (32 random bytes)
 |<---- Encrypted { CHALLENGE: R_B } -----|  [Message 1: Bank -> ATM]
 |                                        |
 |  Compute CARD_PROOF = HMAC(            |
 |    HMAC(K_CARD, ACCOUNT), R_B)         |
 |                                        |
 |---- Encrypted { operation, ACCOUNT, -->|  [Message 2: ATM -> Bank]
 |      amount, CARD_PROOF }              |
 |                                        |  Verify: decrypt -> check HMAC -> check seq#
 |                                        |  -> verify CARD_PROOF -> validate business logic
 |                                        |  -> commit transaction atomically
 |                                        |
 |<---- Encrypted { status, result } -----|  [Message 3: Bank -> ATM]
 |                                        |
 |  Print JSON result, exit(0)            |  Print JSON result
 |---- TCP Close ------------------------>|
\end{verbatim}

For account creation (\texttt{-n}), the flow is identical except
\texttt{CARD\_PROOF} is omitted - the ATM computes and writes the card
secret locally after receiving the success response.

\hypertarget{what-happens-when-things-fail}{%
\subsubsection{What Happens When Things
Fail}\label{what-happens-when-things-fail}}

Every failure has a defined, deterministic outcome. The guiding rules:

\begin{enumerate}
\def\labelenumi{\arabic{enumi}.}
\tightlist
\item
  \textbf{Bank never crashes.} Any bad input -\textgreater{} print
  \texttt{"protocol\_error\textbackslash{}n"}, close connection, resume
  accepting.
\item
  \textbf{ATM exits with a code.} Network/protocol failures
  -\textgreater{} \texttt{exit(63)}. Validation failures or
  bank-reported \texttt{"fail"} -\textgreater{} \texttt{exit(255)}.
  Neither side prints anything to stdout on failure.
\item
  \textbf{No partial state.} The bank commits ledger changes only after
  all validation passes. If anything fails mid-protocol, the ledger is
  untouched.
\item
  \textbf{No information leakage.} The bank gives the same generic
  response regardless of whether a failure was caused by a wrong card,
  missing account, or insufficient funds.
\end{enumerate}

Specific failure points:

\begin{itemize}
\tightlist
\item
  \textbf{TCP connect fails} -\textgreater{} ATM exits 63; bank is
  unaffected.
\item
  \textbf{Decryption/HMAC/sequence/timestamp check fails} (any message,
  either direction) -\textgreater{} receiving side treats it as protocol
  error. Bank prints \texttt{"protocol\_error\textbackslash{}n"} and
  closes. ATM exits 63.
\item
  \textbf{Card proof mismatch} -\textgreater{} Bank closes immediately
  (no response sent), prints
  \texttt{"protocol\_error\textbackslash{}n"}. ATM times out, exits 63.
\item
  \textbf{Business logic failure} (account exists, insufficient funds,
  etc.) -\textgreater{} Bank sends encrypted
  \texttt{\{status:\ "fail"\}}, prints
  \texttt{"protocol\_error\textbackslash{}n"}. ATM exits \texttt{255}.
\item
  \textbf{Response lost in transit} (bank committed but ATM never
  received) -\textgreater{} Bank's ledger is the source of truth; ATM
  exits \texttt{63}. A subsequent balance check reveals the real state.
  This is a fundamental limitation without two-phase commit, and the
  spec accepts it.
\item
  \textbf{3-second timeout} applies to every connection. Prevents
  slowloris, hanging connections, and resource exhaustion.
\item
  \textbf{Message size cap} (\textasciitilde8 KB). Any length prefix
  exceeding this -\textgreater{} immediate close without allocating
  memory.
\end{itemize}

\begin{center}\rule{0.5\linewidth}{0.5pt}\end{center}

\hypertarget{security-model}{%
\subsection{3. Security Model}\label{security-model}}

\hypertarget{attacker-capabilities}{%
\subsubsection{Attacker Capabilities}\label{attacker-capabilities}}

The man-in-the-middle can observe all traffic, modify bytes in transit,
inject new messages, replay old messages, reorder messages, and delay or
drop messages. They have access to the source code but \textbf{not} the
auth file. They may or may not have a card file.

\hypertarget{how-we-defeat-each-attack}{%
\subsubsection{How We Defeat Each
Attack}\label{how-we-defeat-each-attack}}

\begin{longtable}[]{@{}
  >{\raggedright\arraybackslash}p{(\columnwidth - 2\tabcolsep) * \real{0.22}}
  >{\raggedright\arraybackslash}p{(\columnwidth - 2\tabcolsep) * \real{0.78}}@{}}
\toprule\noalign{}
\begin{minipage}[b]{\linewidth}\raggedright
Attack
\end{minipage} & \begin{minipage}[b]{\linewidth}\raggedright
Defence
\end{minipage} \\
\midrule\noalign{}
\endhead
\bottomrule\noalign{}
\endlastfoot
\textbf{Eavesdropping} & AES-256-GCM encrypts all content with K\_ENC.
Without K\_MASTER, nothing is readable. \\
\textbf{Tampering} & GCM's 128-bit auth tag + outer HMAC-SHA256 detect
any modification. \\
\textbf{Replay} & Fresh 32-byte random \texttt{CHALLENGE} per session + monotonic
sequence numbers + timestamp window (±30s). \\
\textbf{Reorder} & Strict sequence number check: expected = previous +
1. \\
\textbf{Injection/Forgery} & Requires \texttt{K\_ENC} and \texttt{K\_MAC} (from auth file)
to produce valid frames. Impossible without the secret. \\
\textbf{Impersonation} & Both sides implicitly authenticate through
shared key derivation. A fake bank cannot produce a valid \texttt{CHALLENGE}; a
fake ATM cannot produce a valid card proof. \\
\textbf{Cross-account access} & Card secret = \texttt{HMAC(K\_CARD,
ACCOUNT\_NAME). Bob's card mathematically cannot authenticate as
Alice}. \\
\textbf{Card forgery} & Requires \texttt{K\_CARD}, which requires the auth
file. \\
\textbf{Timing attacks} & All MAC comparisons use
\texttt{CRYPTO\_memcmp()} (constant-time). \\
\textbf{DoS (connection flood)} & 3-second hard timeout per connection,
thread-per-client handling, and a hard cap of 60 active connections.
Garbage data fails fast at GCM verification
(\textasciitilde microseconds). Bank never crashes. \\
\textbf{Memory exhaustion} & Max message size enforced before buffer
allocation. \\
\end{longtable}

\hypertarget{why-symmetric-only-works}{%
\subsubsection{Why Symmetric-Only
Works}\label{why-symmetric-only-works}}

The auth file eliminates the key-distribution problem. Public-key crypto
(RSA, ECDH, TLS) solves a problem we don't have. AES-256 with a 256-bit
pre-shared key provides 256-bit security - equivalent to a 15,360-bit
RSA key - with a simpler implementation and smaller attack surface.

\begin{center}\rule{0.5\linewidth}{0.5pt}\end{center}

\hypertarget{the-pre-shared-symmetric-key}{%
\subsection{4. The Pre-Shared Symmetric
Key}\label{the-pre-shared-symmetric-key}}

The auth file contains 256 bits of CSPRNG output (the master secret
\texttt{K\_MASTER}). We never use it directly for any operation.
Instead, HKDF-SHA256 derives three domain-separated keys:

\begin{verbatim}
K_MASTER -- HKDF --> K_ENC   (encryption)
                   > K_MAC   (message authentication)
                   > K_CARD  (card secret derivation)
\end{verbatim}

This prevents related-key attacks - even though all keys come from the
same source, they are cryptographically independent. If we used
K\_MASTER directly for both encryption and MAC, a theoretical weakness
in one primitive could compromise the other. HKDF makes this impossible.

\hypertarget{card-file---creation-and-authentication}{%
\subsubsection{Card File - Creation and
Authentication}\label{card-file---creation-and-authentication}}

The card secret is \texttt{HMAC(K\_CARD,\ ACCOUNT\_NAME)} -
deterministic, not random. The bank never stores card secrets; it
recomputes them on the fly from \texttt{K\_CARD} + \texttt{ACCOUNT\_NAME}.

\textbf{Creation (\texttt{-n} mode):}

\begin{verbatim}
ATM                                         BANK
 |                                           |
 |  (no card file exists yet)                |
 |---- Encrypted { create, "bob", 1000 } --> |
 |                                           |  Verify account doesn't exist
 |                                           |  Create account with initial balance
 |<--- Encrypted { success, balance } ------ |
 |                                           |
 |  CARD_SECRET = HMAC(K_CARD, "bob")        |
 |  Write CARD_SECRET to bob.card            |
 |  (or print to stdout with \texttt{-m})     |
 |  (card file never sent to bank)           |
\end{verbatim}

The card secret is computed entirely on the ATM side after a successful
create response. It is never transmitted. If the card file already
exists, the ATM exits 255 before connecting.

\textbf{Authentication (\texttt{-d}, \texttt{-w}, \texttt{-g} modes) -
Challenge-Response:}

\begin{verbatim}
ATM                                         BANK
 |  Read CARD_SECRET from bob.card           |
 |  (or pass card content with \texttt{-C})  |
 |                                           |
 |<--- Encrypted { CHALLENGE: R_B } -------- |  R_B = 32 random bytes (CSPRNG)
 |                                           |
 |  CARD_PROOF = HMAC(CARD_SECRET, R_B)      |
 |                                           |
 |---- Encrypted { op, "bob", amount, -----> |
 |      CARD_PROOF }                         |
 |                                           |  Bank independently computes:
 |                                           |    EXPECTED_SECRET = HMAC(K_CARD, "bob")
 |                                           |    EXPECTED_PROOF  = HMAC(EXPECTED_SECRET, R_B)
 |                                           |  Verify: CARD_PROOF == EXPECTED_PROOF
 |                                           |    (constant-time via CRYPTO_memcmp)
 |                                           |  Match -> proceed; mismatch -> protocol_error
\end{verbatim}

\newpage
\textbf{Why this works:}

\begin{itemize}
\tightlist
\item
  \textbf{The card secret never crosses the wire} - only
  \texttt{CARD\_PROOF\ =\ HMAC(CARD\_SECRET,\ CHALLENGE)} is
  transmitted, and it's a one-time value bound to a fresh random
  challenge.
\item
  \textbf{Uniqueness is mathematical} - HMAC-SHA256 is a PRF; different
  account names yield different secrets with \(2^{128}\) collision
  resistance.
\item
  \textbf{Account binding is intrinsic} - Bob's card yields
  \texttt{HMAC(K\_CARD,\ "bob")} != \texttt{HMAC(K\_CARD,\ "alice")}, so
  Bob's card provably cannot authenticate for Alice's account.
\item
  \textbf{Replay is impossible} - each session uses a fresh 32-byte
  \texttt{R\_B}, so a captured \texttt{CARD\_PROOF} from a previous
  session is useless.
\end{itemize}

\begin{center}\rule{0.5\linewidth}{0.5pt}\end{center}

\hypertarget{design-decisions}{%
\subsection{5. Design Decisions}\label{design-decisions}}

\textbf{AES-256-GCM over AES-CBC + HMAC:} Both are provably secure, but
GCM combines encryption and authentication in one pass, reducing
implementation bugs. It also natively supports Associated Authenticated
Data (AAD), which we use to bind sequence numbers and timestamps to
ciphertexts.

\textbf{Outer HMAC (defence-in-depth):} Strictly redundant given GCM's
auth tag. We keep it as a safety net against hypothetical GCM
implementation bugs and to ensure the unencrypted header fields are
always covered. It costs \textasciitilde200ns per message - negligible.

\textbf{Sequence numbers + timestamps:} For a 3-message protocol,
sequence numbers alone would suffice. We add timestamps (+-30s window)
as cheap insurance against cross-session replay if sequence state is
somehow leaked. If we had to simplify, timestamps would be the first
thing we'd drop.

\textbf{Challenge-response over sending the card secret directly:}
Sending the secret would expose it in the encrypted payload - a problem
if the encryption is ever compromised. The \texttt{CHALLENGE}-response ensures
the card secret never leaves the local machine.

\textbf{Cents as \texttt{uint64} internally:} Avoids all floating-point
precision issues. We only convert to \texttt{double} at JSON output
time. The uint64 range (up to \textasciitilde1.8 x 10\^{}19) is more
than sufficient.

\textbf{C++ with OpenSSL:} OpenSSL is pre-installed on the target
system, requires no internet to build, and provides battle-tested
implementations (AES-NI acceleration, FIPS validation). We use the EVP
API throughout: \texttt{EVP\_aes\_256\_gcm}, \texttt{HMAC()},
\texttt{EVP\_PKEY\_HKDF}, \texttt{RAND\_bytes()},
\texttt{CRYPTO\_memcmp()}, and \texttt{OPENSSL\_cleanse()}.

\begin{center}\rule{0.5\linewidth}{0.5pt}\end{center}

\hypertarget{work-distribution}{%
\subsection{6. Work Distribution}\label{work-distribution}}

\hypertarget{guntas-singh-saran-22110089---crypto-key-management}{%
\subsubsection{Guntas Singh Saran (22110089) - Crypto \& Key
Management}\label{guntas-singh-saran-22110089---crypto-key-management}}

\begin{itemize}
\tightlist
\item
  Auth file generation (CSPRNG) and reading/parsing
\item
  HKDF key derivation (\texttt{K\_ENC}, \texttt{K\_MAC},
  \texttt{K\_CARD})
\item
  AES-256-GCM encrypt/decrypt wrappers (EVP API)
\item
  HMAC-SHA256 wrappers (card secret derivation, card proof, outer MAC)
\item
  RAII \texttt{SecureBuffer} class with \texttt{OPENSSL\_cleanse} in
  destructor
\item
  Constant-time comparison via \texttt{CRYPTO\_memcmp}
\end{itemize}

\textbf{Files:} \texttt{crypto/keys.\{hpp,cpp\}},
\texttt{crypto/aes\_gcm.\{hpp,cpp\}}, \texttt{crypto/hmac.\{hpp,cpp\}},
\texttt{crypto/random.\{hpp,cpp\}}

\hypertarget{jinil-patel-22110184---protocol-networking}{%
\subsubsection{Jinil Patel (22110184) - Protocol \&
Networking}\label{jinil-patel-22110184---protocol-networking}}

\begin{itemize}
\tightlist
\item
  TCP socket layer: server (bind/listen/accept) and client (connect)
\item
  Length-prefixed framing: 4-byte big-endian read/write
\item
  Encrypted session: send = serialize -\textgreater{} encrypt
  -\textgreater{} HMAC -\textgreater{} frame; recv = deframe
  -\textgreater{} verify HMAC -\textgreater{} decrypt -\textgreater{}
  deserialize
\item
  Sequence number tracking (per-direction monotonic counter)
\item
  Timestamp generation and 30s validation
\item
  3-second timeout management via \texttt{setsockopt(SO\_RCVTIMEO)}
\item
  Message size cap enforcement
\item
  SIGTERM signal handler (bank graceful shutdown)
\end{itemize}

\textbf{Files:} \texttt{protocol/framing.\{hpp,cpp\}},
\texttt{protocol/session.\{hpp,cpp\}}

\hypertarget{bhavik-patel-22110047---bank-server-business-logic}{%
\subsubsection{Bhavik Patel (22110047) - ATM Client (Business Logic)
\& Bank Server + UI Framework (Business
Logic)}\label{bhavik-patel-22110047---bank-server-business-logic}}

\begin{itemize}
\tightlist
\item
  Bank \texttt{main.cpp}: CLI parsing, startup sequence, auth file
  creation
\item
  Account ledger:
  \texttt{std::unordered\_map\textless{}std::string,\ uint64\_t\textgreater{}}
  protected by \texttt{std::mutex}
\item
  Thread-per-client handling with connection cap
  (\texttt{MAX\_ACTIVE\_CLIENTS = 60})
\item
  Per-connection handler: decrypt request -\textgreater{} verify
  CARD\_PROOF -\textgreater{} dispatch operation
\item
  Operations: create account (check doesn't exist, balance \textgreater=
  10.00), deposit (\textgreater{} 0), withdraw (sufficient funds),
  get-balance
\item
  Atomic commit: ledger modified only after all validations pass
\item
  JSON output to stdout (bank side), flush after every line
\item
  All error paths -\textgreater{}
  \texttt{"protocol\_error\textbackslash{}n"}, close connection,
  continue serving
\item
  ATM \texttt{main.cpp}: CLI parsing (POSIX/getopt), mode dispatch
  (-n/-d/-w/-g, \texttt{-m}, \texttt{-C}, \texttt{--help})
\item
  Card file I/O: create (write HMAC secret or \texttt{-m} to stdout),
  read (parse secret or \texttt{-C} direct content)
\item
  Transaction execution: connect -\textgreater{} receive \texttt{CHALLENGE}
  -\textgreater{} compute CARD\_PROOF -\textgreater{} send request
  -\textgreater{} receive response -\textgreater{} print JSON
\item Laid the groundwork for the UI integration (see \texttt{ui/backend/}) - a local Node.js/Express server that spawns \texttt{atm} processes on demand, never exposing secrets to the browser.
\end{itemize}

\textbf{Files:} \texttt{build/src/atm/main.cpp},\texttt{build/src/bank/main.cpp}, \texttt{build/include/common/exitcodes.hpp},
\texttt{build/include/common/types.hpp}, \texttt{ui/backend/}, \texttt{ui/frontend/}

\hypertarget{ui-frontend-backend}{%
\subsubsection{UI Frontend \& Local Backend (Web UI)}}

\begin{itemize}
\tightlist
\item
  React + Vite SPA in \texttt{ui/frontend} (hosted publicly)
\item
  Node.js/Express backend in \texttt{ui/backend} (runs locally on each user machine)
\item
  Backend spawns \texttt{atm} with argv, never exposing auth or card secrets to the browser
\item
  Deployment and runtime model documented in \texttt{ui/README.md} and \texttt{ui/DEPLOYMENT.md}
\end{itemize}

\newpage
\begin{verbatim}
Browser (Public SPA)
      |
      | HTTPS (static assets)
      v
Public Frontend Host (GitHub Pages)
      |
      | fetch http://127.0.0.1:4000/api/*
      v
Local Node.js Backend (user machine)
      |
      | spawn(atm, argv)
      v
Local atm binary + local bank.auth + card files
      |
      | TCP
      v
Remote bank server (public IP:port)
\end{verbatim}

\textbf{Files:} \texttt{ui/frontend/*}, \texttt{ui/backend/*},
\texttt{ui/README.md}, \texttt{ui/DEPLOYMENT.md}, \texttt{.github/workflows/ui-gh-pages.yml}

\hypertarget{rishabh-jogani-23110276---atm-client-business-logic-input-validation}{%
\subsubsection{Rishabh Jogani (23110276) - Input
Validation}\label{rishabh-jogani-23110276---atm-client-business-logic-input-validation}}

\begin{itemize}
\tightlist
\item
  Full input validation: account name regex
  \texttt{{[}\_\textbackslash{}-\textbackslash{}.0-9a-z{]}\{1,122\}},
  amount format
  \texttt{(0\textbar{}{[}1-9{]}{[}0-9{]}*)\textbackslash{}.{[}0-9{]}\{2\}},
  file names
  \texttt{{[}\_\textbackslash{}-\textbackslash{}.0-9a-z{]}\{1,127\}}
  (not \texttt{.}/\texttt{..}), IPv4 dotted-decimal, port 1024--65535,
  duplicate flag detection
\item
  JSON output formatting: cents -\textgreater{} correct decimal
  representation, proper precision
\item
  CMakeLists.txt, wrapper Makefile, build system integration
\item
  Exit code logic: 0 (success), 63 (protocol/network error), 255
  (validation/business error)
\end{itemize}

\textbf{Files:} \texttt{build/src/validator/validator.cpp},
\texttt{build/include/validator/validator.hpp},
\texttt{build/CMakeLists.txt}, \texttt{build/Makefile}

\hypertarget{integration-points}{%
\subsubsection{Integration Points}\label{integration-points}}

Members A and B produce the shared \texttt{common} static library
(crypto + protocol + validators). Members C and D build the
\texttt{bank} and \texttt{atm} executables that link against it. The
interface boundary is clean: Member C and D call \texttt{session.send()}
/ \texttt{session.recv()} and get back decrypted, verified JSON - they
never touch raw crypto.

\begin{center}\rule{0.5\linewidth}{0.5pt}\end{center}

\hypertarget{build-system}{%
\subsection{7. Build System}\label{build-system}}

We use CMake with a thin Makefile wrapper so that
\texttt{cd\ build\ \&\&\ make} works as the spec requires:

\newpage
\begin{verbatim}
secure-atm/
|-- build/
|   |-- Makefile            # Build wrapper (run `make` here)
|   |-- CMakeLists.txt      # CMake build configuration
|   |-- README.md           # Build/run usage for CLI
|   |-- include/
|   |   |-- common/{types,exitcodes}.hpp
|   |   |-- crypto/{random,keys,aes_gcm,hmac}.hpp
|   |   |-- protocol/{framing,session}.hpp
|   |   |-- validator/validator.hpp
|   |   |-- nlohmann/json.hpp
|   |-- src/
|   |   |-- crypto/{random,keys,aes_gcm,hmac}.cpp
|   |   |-- protocol/{framing,session}.cpp
|   |   |-- validator/validator.cpp
|   |   |-- bank/main.cpp
|   |   |-- atm/main.cpp
|   |-- tests/
|       |-- test_framework.hpp
|       |-- test_validator.cpp
|       |-- test_crypto.cpp
|       |-- test_protocol.cpp
|       |-- test_integration.sh
|       |-- dos_attack.cpp
|-- ui/
    |-- frontend/           # React + Vite SPA (public hosting)
    |-- backend/            # Local Node.js backend (spawns atm)
\end{verbatim}

Dependencies: C++17 compiler (g++), OpenSSL (\texttt{libssl-dev}),
pthreads - all available offline on the target system. No internet
required.

\begin{center}\rule{0.5\linewidth}{0.5pt}\end{center}

\hypertarget{summary}{%
\subsection{8. Summary}\label{summary}}

The design uses symmetric-only crypto anchored on a pre-shared 256-bit
secret. HKDF derives domain-separated keys. AES-256-GCM provides
confidentiality and integrity. An HMAC-based challenge-response
authenticates card possession without sending secrets over the wire.
Sequence numbers, timestamps, and a fresh random \texttt{CHALLENGE} per session
prevent replay and reordering. The bank never crashes, commits
atomically, and times out all connections at 3 seconds. A web UI is
provided via a public frontend with a local backend that spawns \texttt{atm}.
The whole thing compiles with \texttt{make} from source on an offline machine.

\end{document}
